% !TeX root = ../main.tex
\section{Заключение}

В ходе выполнения лабораторной работы были изучены архитектура сетевых СРНС, принципы взаимодействия их базовых сегментов и математические основы псевдодальномерного метода позиционирования.

На основе натурных измерений с использованием бытовой НАП (смартфона) была экспериментально исследована зависимость точности навигационных определений от окружающей радиотехнической обстановки:
\begin{itemize}
    \item \textbf{Открытая местность:} погрешность составила $2\text{ м}$, что полностью согласуется с паспортными характеристиками автономного кодового режима SPP ($2\dots5\text{ м}$) при благоприятной геометрической конфигурации созвездия (низкие значения DOP) и минимальном уровне переотражений.
    \item \textbf{Малоэтажная застройка:} погрешность возросла до $7\text{ м}$ вследствие частичного затенения горизонта и появления интерференционных задержек от переотражения радиоволн.
    \item \textbf{Плотная высотная застройка (9-этажное здание):} зафиксировано максимальное линейное отклонение $37\text{ м}$. Резкая деградация точности обусловлена классическим эффектом «городского каньона», сочетающим интенсивную многолучевость и значительное ухудшение геометрического фактора (рост PDOP) из-за блокировки прямого сигнала от большинства космических аппаратов.
\end{itemize}

Эксперимент наглядно подтвердил, что в реальных городских условиях ключевым дестабилизирующим фактором для бытовой некорректируемой НАП являются переотражения от окружающих построек и ограничение рабочего сектора видимости спутников.
